% docs/physics/latex/05_delta.tex - SIMC++ physics note 05: (e,e'p)pi in the Delta region.
% Standalone PDF docs/physics/pdf/SIMCpp_Physics_05_delta.pdf (copied to src/reactions/delta/ and src/physics/delta/);
% a chapter of the book SIMC++ Physics Reference. Rules: docs/physics/README.md; structure: template_reaction.tex.
\documentclass{simcnote}
\simccommon
\begin{document}
\simcnotemeta{05}{(e,e\pr p)\texorpdfstring{$\pi$}{pi}, \texorpdfstring{$\Delta$}{Delta} region}{Proton detection in the \texorpdfstring{$\Delta$}{Delta} region: H, D, A(e,e\pr p)\texorpdfstring{$\pi$}{pi}}%
  {Reaction plug-in \texttt{delta\_production} with the physics model \texttt{delta}: pion electroproduction with the
   proton detected and the pion missing; the weight is phase space times virtual-photon flux}
\simcnotemaketitle
\simcabstract{Pion electroproduction in the $\Delta(1232)$ region seen through the recoil proton: the electron and the
proton are detected, the pion is the missing particle. Generation and kinematics are those of exclusive pion production
(note 04) with the roles exchanged: the detected hadron is the proton, the recoil is a pion, and one of the two proton
energies allowed by the kinematics is chosen at random. The physics model computes the photon--nucleon centre-of-mass
frame, the jacobian and the virtual-photon flux, but the event weight contains no cross section: the cross-section line
of the code is commented out. Notation, glossary and common physics: note 00.}
\simcnotetoc

% =====================================================================================================================
\section{The reaction in SIMC++}\lbl{sec:reaction}

\subsection{Process and scope}
The reaction is $e + p \to e' + p + \pi$ with the scattered electron and the proton detected and the pion
reconstructed as the missing mass; the code comments describe it as ``(e,e\pr p)pi'' and as a first version with a
``shortcut'' cross section, set up for a proton target. The struck nucleon is a proton ($M_N=M_p$), the detected hadron
is a proton ($M_h=M_p$) and the recoil is a pion of mass $M_X=M_\pi=139.57018$ MeV, the charged-pion mass.

\begin{simcwarning}
The event weight of this reaction contains \emph{no cross section}. The model returns the photon flux times the
centre-of-mass jacobian times 1 (\secr{sec:xsec}); the line that multiplied by the cross section is commented out in
the code (``Need new xsec model.''). Yields, \code{wtcontr} and the column \code{Weight} of a Delta run are therefore
phase-space distributions weighted by the virtual-photon flux, not rates. SIMC++ keeps this in both modes; whether a
real cross section should be provided is an open question (\secr{sec:origin}).
\end{simcwarning}

What is simulated: the generation of note 04 (electron energy, both slope pairs, Fermi motion for D and $A\ge3$), the
two-body kinematics with a pion recoil, the photon--nucleon centre-of-mass frame and its jacobian for the proton, the
event chain of note 00 (radiation on the three legs, energy loss and multiple scattering of the proton, Coulomb
correction). What is not simulated: a $\Delta$ line shape or any resonance dynamics (no cross section at all), the
$n\pi^+$ channel, final-state interactions. The Fortran stops every deuterium or heavier-target run before the first
event; SIMC++ runs them in the default mode only (\secr{sec:origin}), with no Fortran reference for those cases.

\subsection{Components}
\begin{simctable}{footnotesize}{L{0.15\linewidth}L{0.16\linewidth}L{0.15\linewidth}L{0.25\linewidth}L{0.29\linewidth}}
\caption{The SIMC++ components of (e,e\pr p)$\pi$ (one component per row).}\lbl{tab:components}\\
\toprule\rowcolor{simchead}\simcth{role} & \simcth{interface} & \simcth{registered name} & \simcth{class} & \simcth{file}\\\midrule\endfirsthead
\toprule\rowcolor{simchead}\simcth{role} & \simcth{interface} & \simcth{registered name} & \simcth{class} & \simcth{file}\\\midrule\endhead
\bottomrule\endlastfoot
reaction & \code{IReaction} & \code{delta\_production} & \code{reactions::DeltaProduction} & \file{reactions/delta/DeltaProduction.cpp}\\
generation and completion & (base class) & -- & \code{reactions::PionProduction} & \file{reactions/pion/PionProduction.cpp}\\
reconstruction & (base class) & -- & \code{reactions::EepReactionBase} & \file{reactions/EepReactionBase.cpp}\\
centre-of-mass transform & (function) & -- & \code{reactions::transformToCm} & \file{reactions/EventChain.cpp}\\
weight model & \code{IPhysicsModel} & \code{delta} & \code{delta::DeltaModel} & \file{physics/delta/DeltaModel.cpp}\\
model kernel & (function) & -- & \code{delta::peedelta} & \file{physics/delta/DeltaModel.cpp}\\
Blok fit (tree only) & (function) & -- & \code{pion::sigBlok} & \file{physics/pion/PionCrossSections.cpp}\\
missing energy, $A\ge3$ & (class) & -- & \code{nuclear::SfTable} & \file{physics/nuclear/SpectralFunctionTables.cpp}\\
beam, target, radiation, arms, outputs & as in note 00 & \code{standard}, \code{hallc\_cryo}, \code{simc\_rad}, \ldots & & \ctabr{tab:plugins}\\
\end{simctable}

\subsection{Selection and override}
The input \code{doing\_delta} = 1 selects \code{delta\_production} with the model \code{delta} (\ctabr{tab:defaults}).
The case follows the target: H for $\mathrm{nint}(A)=1$, D for 2, $A\ge3$ otherwise; for $A\ge2$ a warning says that the
model is set up for a proton target only. The set-up uses $M_h=M_N=M_p$, $M_X=M_\pi$, $M_A=M_p$ for H and $M_A=M_d$,
$M_\mathrm{rec}=M_n$ for D; the missing-mass cut window is opened to $\pm10^6$ MeV. The input \code{which\_pion} is also
read by this reaction: with 2 or 3 the recoil pion is replaced by a $\Delta$ mass drawn as in note 04, and with 1, 3 or
11 the $\pi^-$ branch of the Blok fit is used for the tree column \code{sigcm}; keep \code{which\_pion} = 0.

\begin{simctable}{footnotesize}{L{0.3\linewidth}L{0.38\linewidth}L{0.32\linewidth}}
\caption{JSON keys of this reaction and its model (one key per row).}\lbl{tab:json}\\
\toprule\rowcolor{simchead}\simcth{key} & \simcth{effect} & \simcth{example}\\\midrule\endfirsthead
\toprule\rowcolor{simchead}\simcth{key} & \simcth{effect} & \simcth{example}\\\midrule\endhead
\bottomrule\endlastfoot
\code{pipeline.reaction.data\_dir} & directory of \code{delta\_relativistic\_bw.inp} (read only for \code{which\_pion} 2/3) & \code{"data/pion"}\\
\code{pipeline.reaction.verbose} & print the Fortran warnings of the reaction & \code{true}\\
\code{pipeline.physics.verbose} & print the warnings of the transform & \code{true}\\
\end{simctable}
Because the reaction only asks the model for one number, a real cross section can be plugged in without touching the
reaction: a model registered under a new name, selected with
\code{\{"pipeline": \{"reaction": "delta\_production", "physics": "my\_delta"\}\}}, must return
$\mathrm{d}^5\sigma/\mathrm{d}E'\mathrm{d}\Omega_e\mathrm{d}\Omega_p$ in $\mu$b/(MeV sr$^2$) with the proton solid angle in
the laboratory (\secr{sec:require}).

\subsection{Event flow}
\begin{simctable}{footnotesize}{L{0.05\linewidth}L{0.33\linewidth}L{0.62\linewidth}}
\caption{One try of (e,e\pr p)$\pi$ (one step per row; the general chain is \ctabr{tab:chain}).}\lbl{tab:flow}\\
\toprule\rowcolor{simchead}\simcth{\#} & \simcth{C++ call} & \simcth{what is specific to this reaction}\\\midrule\endfirsthead
\toprule\rowcolor{simchead}\simcth{\#} & \simcth{C++ call} & \simcth{what is specific to this reaction}\\\midrule\endhead
\bottomrule\endlastfoot
1 & \code{StandardBeam::generate} & nothing\\
2 & \code{PionProduction::generate} & both slope pairs, $E'$; Fermi momentum and $E_m$ for D/$A\ge3$ (default mode only)\\
3 & \code{PionProduction::completeEv} & the proton energy with a random choice of root (\eqr{eq:roots}), $W$, $\varepsilon$, $-t$, $\varphi_{pq}$\\
4 & \code{SimcRadiator::radiate} & tail-1 limits (\secr{sec:rad}); a tail-1 photon draws a new root\\
5 & \code{MonteCarlo::run} & the proton crosses the hadron arm\\
6 & \code{EepReactionBase::completeRecon} & missing mass of the pion (\eqr{eq:recmm})\\
7 & \code{DeltaProduction::completeMain} & $S=1$, \code{DeltaModel::sigccMain} (\eqr{eq:peedelta}), $n_\mathrm{tgt}=1$\\
\end{simctable}

% =====================================================================================================================
\section{Kinematics, generated variables and generation volume}\lbl{sec:gen}
The reaction uses the generation of exclusive pion production unchanged; note 04, \S2 has the formulas. The generated
variables are the electron slopes $y'_e$, $x'_e$ and the proton slopes $y'_h$, $x'_h$ (flat in the edges), the electron
energy $E'$ (flat in the range narrowed by the summed-energy limits, with the generation weight), and for D and $A\ge3$
the Fermi momentum $\vec p_F$, the missing energy $E_m$ and the struck-nucleon energy $E_F$:
\begin{feqa}{reactions/pion/PionProduction.cpp: PionProduction::generate}{event.f:276-286, 300-315, 336-372}
E' &= E'^\mathrm{lo}+\xi\,(E'^\mathrm{hi}-E'^\mathrm{lo}),\qquad [E'^\mathrm{lo},E'^\mathrm{hi}]=\bigl[\max(E'^\mathrm{gen}_\mathrm{min},\Sigma_\mathrm{min}),\ \min(E'^\mathrm{gen}_\mathrm{max},\Sigma_\mathrm{max})\bigr],\nonumber\\
w_\mathrm{gen} &\leftarrow w_\mathrm{gen}\,\frac{E'^\mathrm{hi}-E'^\mathrm{lo}}{E'^\mathrm{gen}_\mathrm{max}-E'^\mathrm{gen}_\mathrm{min}},\nonumber\\
E_F &= M_A-\sqrt{(M_A-M_N+E_m)^2+p_F^2}\ \ (\text{D},\ A\ge3),\qquad E_F=M_N\ \ (\text{H}),
\lbl{eq:gen}
\end{feqa}
the Fermi momentum coming from the momentum distribution and, for $A\ge3$, $E_m$ from the spectral function (note 04,
\S2.2). The limits are those of note 04 with $M_X=M_\pi$ and $M_h=M_p$:
\begin{feqa}{setup/SetupLimits.cpp: BuildLimits}{init.f:411-413, 456-469, 477-480}
E_\gamma^\mathrm{tot,max} &= E_b^\mathrm{max}+M_p-M_\pi-E'_\mathrm{min}-E_{h,\mathrm{min}}-E_m^\mathrm{min}-T_\mathrm{rec}^\mathrm{min},\nonumber\\
\Sigma_\mathrm{max} &= \min\bigl(E_b^\mathrm{max}+M_p-M_\pi-E_{h,\mathrm{min}}-E_m^\mathrm{min}-T_\mathrm{rec}^\mathrm{min},\ E'_\mathrm{max}+E_{\gamma2}^\mathrm{max}\bigr)+d_E,
\lbl{eq:limits}
\end{feqa}
and correspondingly for $\Sigma_\mathrm{min}$ and the vertex limits ($E_{h}$ is the proton energy here). The generation
volume is
\begin{feq}{reactions/EepReactionBase.cpp: EepReactionBase::genVolume}{simc.f:383-397}
V_\mathrm{gen} = \Delta\Omega_e\,\Delta\Omega_p\,\bigl(E'^\mathrm{gen}_\mathrm{max}-E'^\mathrm{gen}_\mathrm{min}\bigr)\quad[\mathrm{sr^2\,MeV}].
\lbl{eq:genvol}
\end{feq}

% =====================================================================================================================
\section{Completing the event}\lbl{sec:complete}
\subsection{The proton energy}
The proton energy solves the quadratic of exclusive pion production with $M_h=M_p$ and the pion as recoil; unlike the
pion reaction, one random number chooses between the two roots:
\begin{feqa}{reactions/pion/PionProduction.cpp: PionProduction::completeEv}{event.f:634-698}
a &= -|\vec q|\,(\hat q\cdot\hat p) - p_F\,(\hat p_F\cdot\hat p),\quad
b = |\vec q|^2+p_F^2+2|\vec q|p_F(\hat p_F\cdot\hat q),\quad c=\nu+E_F\ \ (\nu+M_A\text{ for H}),\nonumber\\
\tau &= c^2-b+M_p^2-M_\pi^2,\quad Q_A=4(a^2-c^2),\quad Q_B=4c\tau,\quad Q_C=-4a^2M_p^2-\tau^2,\nonumber\\
E_\pm &= \frac{-Q_B\mp\sqrt{Q_B^2-4Q_AQ_C}}{2Q_A},\qquad
E_p = \begin{cases}E_+ & \xi\le0.5\\ E_- & \xi>0.5\end{cases},\qquad P_p=\sqrt{E_p^2-M_p^2}.
\lbl{eq:roots}
\end{feqa}
The try fails for a negative discriminant, when $E_+<0$ (tested before the choice), when the pion energy $c-E_p\le M_\pi$
or when $E_p\le M_p$. When both roots are physical each is generated in half of the tries with the full weight; no factor
2 enters $w_\mathrm{gen}$ or $J$ (\secr{sec:origin}). The rest of the completion is that of note 04: $W$, $\varepsilon$,
$\theta_{pq}$, $\varphi_{pq}$, and
\begin{feq}{reactions/pion/PionProduction.cpp: PionProduction::pionMainKinematics}{event.f:728-731}
-t = Q^2-M_p^2+2\nu E_p-2P_p|\vec q|\cos\theta_{pq},
\lbl{eq:t}
\end{feq}
the momentum transfer to the \emph{proton}, $-(q-p_p)^2$, under the name of the pion $t$. $T_\mathrm{rec}=0$ for H.
With \code{which\_pion} 2 or 3 the recoil mass $M_\pi$ is replaced by a sampled $\Delta$ mass (note 04, \S3.1), which lies
far above the available energy in the $\Delta$ region: every completion then fails.

\subsection{Reconstruction}\lbl{sec:recon}
\begin{feqa}{reactions/EepReactionBase.cpp: EepReactionBase::completeRecon}{event.f:1325-1335, 1353-1354}
E_m^\mathrm{rec} &= \nu^\mathrm{rec}+M_p-E_p^\mathrm{rec},\qquad
m_N = \operatorname{sign}(m_N^2)\sqrt{|m_N^2|},\quad m_N^2=(E_m^\mathrm{rec})^2-(P_m^\mathrm{rec})^2,\nonumber\\
m_A^2 &= (\nu^\mathrm{rec}+M_A-E_p^\mathrm{rec})^2-(P_m^\mathrm{rec})^2,\qquad
-t^\mathrm{rec} = (Q^\mathrm{rec})^2-M_p^2+2\bigl(\nu^\mathrm{rec}E_p^\mathrm{rec}-P_p^\mathrm{rec}|\vec q^\mathrm{rec}|\cos\theta_{pq}^\mathrm{rec}\bigr),
\lbl{eq:recmm}
\end{feqa}
so that \code{missmass} ($m_N$) is the reconstructed mass of the missing pion, peaking at $M_\pi$ for hydrogen without
radiation.

% =====================================================================================================================
\section{Jacobians}\lbl{sec:jac}
The event jacobian is that of the four generated slopes,
\begin{feq}{reactions/pion/PionProduction.cpp: PionProduction::completeEv}{event.f:1015-1025}
J = \bigl(1+x_e'^2+y_e'^2\bigr)^{-3/2}\bigl(1+x_p'^2+y_p'^2\bigr)^{-3/2},
\lbl{eq:jac}
\end{feq}
and the model applies the jacobian $J_\mathrm{cm}=|\partial(t,\varphi^*)/\partial(\cos\theta_{pq},\varphi_{pq})|$ of the
photon--nucleon centre-of-mass transformation of note 04 (\S4.2), with the proton in the place of the pion: $\theta^*$,
$\varphi^*$ are the centre-of-mass angles of the proton relative to $\vec q^{\,*}$, and $t$ is that of \eqr{eq:t}.

% =====================================================================================================================
\section{Cross-section model}\lbl{sec:xsec}
\subsection{What the reaction requires from a model}\lbl{sec:require}
\code{DeltaProduction::completeMain} sets $S=1$, calls \code{sigccMain(vertex, main)} of the run's
\code{IPhysicsModel}, multiplies by the Coulomb factor and uses $n_\mathrm{tgt}=1$ for every target. The model reads the
same quantities as the pion model (the vertex record, $\varepsilon$, $-t$, $\varphi_{pq}$, the Fermi momentum, $M_N$) and
must return $\mathrm{d}^5\sigma/\mathrm{d}E'\mathrm{d}\Omega_e\mathrm{d}\Omega_p$ in $\mu$b/(MeV sr$^2$); \code{sigccRecon}
returns 1.

\subsection{The model \texttt{delta}}
With $s$, $W^*$, $k_\mathrm{eq}$, $\Gamma$ and $f$ of note 04 (\S5.2) and $J_\mathrm{cm}$ of \secr{sec:jac}:
\begin{feqa}{physics/delta/DeltaModel.cpp: peedelta}{physics\_delta.f:50-132}
s &= (\nu+E_F)^2-\bigl(|\vec q|+p_F\cos\theta_F'\bigr)^2-\bigl(p_F\sin\theta_F'\bigr)^2,\qquad
k_\mathrm{eq} = \frac{s-M_p^2}{2M_p},\qquad
\Gamma = \frac{\alpha}{2\pi^2}\frac{E'}{E_\mathrm{in}}\frac{k_\mathrm{eq}}{Q^2}\frac{1}{1-\varepsilon},\nonumber\\
f &= \frac{1}{1-p_F\hat p_{F,z}/E_F}\frac{M_p}{E_F},\qquad
\sigma = 1\cdot J_\mathrm{cm}\cdot\Gamma\cdot f,
\lbl{eq:peedelta}
\end{feqa}
that is, the weight of a constant centre-of-mass cross section $\mathrm{d}^2\sigma/\mathrm{d}t\,\mathrm{d}\varphi^*=1$
$\mu$b/(MeV$^2$ rad). The commented-out line multiplied by the Blok fit below, a $p(e,e'\pi^+)n$ fit, which would not be a
$\Delta\to p\pi$ cross section either. The model still evaluates that fit and writes it to the tree column \code{sigcm}:
\begin{feqa}{physics/pion/PionCrossSections.cpp: sigBlok}{physics\_pion.f:427-462}
\sigma_L &= 27.8\,e^{-11.5|t|},\quad \sigma_T = 50|t|\,e^{-5|t|},\quad \sigma_{LT}=0,\quad
\sigma_{TT} = -(4\sigma_L+0.5\sigma_T)\sin^2\theta^*,\nonumber\\
\pi^-:&\quad \sigma_T,\ \sigma_{TT}\ \leftarrow\ 0.25\,(1+3e^{-10|t|})\,\sigma_T,\ \sigma_{TT},\nonumber\\
F_\pi &= \frac{1}{1+1.65Q^2+0.5Q^4},\quad
\sigma_L\leftarrow\sigma_L\frac{Q^2F_\pi^2}{0.1215},\quad \sigma_T\leftarrow\frac{\sigma_T}{0.3+Q^2},\quad
\sigma_{TT}\leftarrow\frac{\sigma_{TT}}{0.3+Q^2},\nonumber\\
\frac{\mathrm{d}^2\sigma}{\mathrm{d}t\,\mathrm{d}\varphi^*}\Big|_\mathrm{Blok} &=
\bigl(\sigma_T+\varepsilon\sigma_L+\varepsilon\cos2\varphi^*\sigma_{TT}+\sqrt{2\varepsilon(1+\varepsilon)}\cos\varphi^*\sigma_{LT}\bigr)
\frac{15.333}{(s-M_N^2)^2}\,\frac{1}{2\pi\cdot10^6},
\lbl{eq:sigblok}
\end{feqa}
with $Q^2$, $|t|$, $s$, $M_N$ in GeV units, the result in $\mu$b/(MeV$^2$ rad); the $\pi^-$ branch is used for
\code{which\_pion} 1, 3, 11. The code comments: the fit reproduces the data of Brauel et al.\ \citer{brauel} at $W=2.19$
GeV and $Q^2=0.70$ GeV$^2$ (``MODIFIED BY HAND''), 0.1215 is $Q^2F_\pi^2$ at $Q^2=0.7$, and 15.333 is $(W^2-M_p^2)^2$ at
$W=2.19$ GeV (with $M_p$ in GeV this gives 15.337); the coefficients are transcribed from the code and were not compared
with the paper. Here $t$ and the angles are those of the proton. In the Fortran the arguments $t$, $Q^2$, $s$ are passed
with the wrong floating-point kind, so that \code{sigcm} holds meaningless numbers; the compat mode reproduces them, the
default mode passes the intended values.

\subsection{Weight and units}
\begin{feq}{reactions/delta/DeltaProduction.cpp: DeltaProduction::completeMain}{event.f:1514-1516, 1536-1538, 1561-1562}
w = J\cdot w_\mathrm{gen}\cdot\sigma\quad(n_\mathrm{tgt}=1),\qquad \sigma\leftarrow\sigma\,(1+V_C/E_b)^2\ \text{with the Coulomb correction},
\lbl{eq:weight}
\end{feq}
in the units of \eqr{eq:peedelta} with the unit cross section. The \texttt{.hist} line \code{CENTRAL.sigcc} is the model
at the central kinematics, with $\varepsilon=-t=\varphi_{pq}=0$ in the compat mode and the central values in the
default mode (\code{PionProduction::prepareCentral}).

% =====================================================================================================================
\section{Radiation, energy loss and decays}\lbl{sec:rad}
The three legs radiate, the proton with $\lambda_3$ of \ceqr{eq:lambda}. Tails 2 and 3 have the limits of note 04
(\S6.1) with the proton energy and edges. For tail 1 the Fortran sets no limits for this reaction (the case list of
the incident-electron tail omits it): the limits keep the value 0, the photon generation returns a zero weight and
every try that chooses tail 1 fails, which the compat mode reproduces. The default mode uses the pion case:
\begin{feq}{radiation/Radc.cpp: generateRad}{radc.f:297-308}
E_{\gamma1}\in\bigl[-d_E,\ \min(\Sigma_\mathrm{max}-E',\ E_{\gamma1}^\mathrm{max})+d_E\bigr]\quad(\text{default mode}).
\lbl{eq:tail1}
\end{feq}
After a tail-1 photon the completion runs again and draws a new root choice. The proton loses energy and scatters like
any hadron (\csecr{sec:eloss}). Nothing decays: SIMC++ refuses \code{doing\_decay} for a proton.

% =====================================================================================================================
\section{Inputs and outputs}\lbl{sec:io}
\subsection{Inputs}
\begin{simctable}{footnotesize}{L{0.24\linewidth}L{0.12\linewidth}L{0.08\linewidth}L{0.56\linewidth}}
\caption{Inputs specific to (e,e\pr p)$\pi$ (one input per row; the common inputs are in note 01, \S7, and the user
manual, chapter 4).}\lbl{tab:inputs}\\
\toprule\rowcolor{simchead}\simcth{input} & \simcth{unit} & \simcth{default} & \simcth{meaning here}\\\midrule\endfirsthead
\toprule\rowcolor{simchead}\simcth{input} & \simcth{unit} & \simcth{default} & \simcth{meaning here}\\\midrule\endhead
\bottomrule\endlastfoot
\code{doing\_delta} & flag & 0 & 1 selects this reaction\\
\code{which\_pion} & code & 0 & keep 0 (2/3: $\Delta$ recoil mass, every try fails; 1/3/11: $\pi^-$ branch of \code{sigcm})\\
\code{targ\%A} & -- & 0 & H, D or $A\ge3$ (D and $A\ge3$ run in the default mode only)\\
\code{targ\%Z} & -- & 0 & Coulomb correction (not $n_\mathrm{tgt}$)\\
\code{hadron\_arm} & code & 0 & proton spectrometer\\
\code{spec\%p\%P} & MeV/$c$ & 0 & $P_h^0$, the proton central momentum\\
\code{spec\%p\%theta} & deg & 0 & $\theta_h^0$\\
\code{SPedge\%p\%yptar\%min} & mrad & 0 & generation range of $y'_p$, lower limit\\
\code{SPedge\%p\%yptar\%max} & mrad & 0 & generation range of $y'_p$, upper limit\\
\code{SPedge\%p\%xptar\%min} & mrad & 0 & generation range of $x'_p$, lower limit\\
\code{SPedge\%p\%xptar\%max} & mrad & 0 & generation range of $x'_p$, upper limit\\
\code{SPedge\%p\%delta\%min} & \% & 0 & proton energy edges (limits, tail 3)\\
\code{SPedge\%p\%delta\%max} & \% & 0 & proton energy edges (limits, tail 3)\\
\end{simctable}

\subsection{Outputs}
% docs/physics/latex/frag_pion_columns.tex - the meson tree layout of the pion and Delta reactions, \input by the notes
% 04 (pion_production) and 05 (delta_production). C++: output/NtupleLayout.cpp (UpstreamMesonColumns,
% FillUpstreamNtuple; legacy layout PionNtupleColumns, FillPionNtuple). Fortran: NtupleInit.f:34-224 and
% results_write.f:120-286 (upstream 59ccdf6). One column per row. No displayed equations here.
The sink \code{root\_tree} writes one entry of the tree \code{h10} per success. With the layout of upstream SIMC 59ccdf6
(the default) the tree has the 64 columns of \tabr{tab:columns}; \code{-{}-tree-layout legacy} gives the layout of SIMC
294f662. Columns 1--24 are those of the (e,e\pr p) layout: \code{hs} is the arm on the beam right, \code{ss} the arm on
the beam left, whatever particle each detects; the reconstructed columns come from the \code{recon} record, the
``true'' ones (suffix \code{i}) from the \code{orig} record for the arm columns and from the \code{vertex} record
for the physics columns. ``Hadron'' is the detected hadron: the pion in note 04, the proton in note 05.

\begin{simctable}{footnotesize}{L{0.06\linewidth}L{0.15\linewidth}L{0.13\linewidth}L{0.66\linewidth}}
\caption{Tree columns of the pion and Delta reactions (default layout; one column per row).}\lbl{tab:columns}\\
\toprule\rowcolor{simchead}\simcth{\#} & \simcth{name} & \simcth{unit} & \simcth{content}\\\midrule\endfirsthead
\toprule\rowcolor{simchead}\simcth{\#} & \simcth{name} & \simcth{unit} & \simcth{content}\\\midrule\endhead
\bottomrule\endlastfoot
1 & \code{hsdelta} & \% & reconstructed $\delta$, \code{hs} arm\\
2 & \code{hsyptar} & rad & reconstructed $y'$, \code{hs} arm\\
3 & \code{hsxptar} & rad & reconstructed $x'$, \code{hs} arm\\
4 & \code{hsytar} & cm & reconstructed horizontal target position, \code{hs} arm\\
5 & \code{hsxfp} & cm & focal-plane $x$, \code{hs} arm\\
6 & \code{hsxpfp} & rad & focal-plane $x'$, \code{hs} arm\\
7 & \code{hsyfp} & cm & focal-plane $y$, \code{hs} arm\\
8 & \code{hsypfp} & rad & focal-plane $y'$, \code{hs} arm\\
9 & \code{hsdeltai} & \% & true $\delta$ after radiation, \code{hs} arm\\
10 & \code{hsyptari} & rad & true $y'$, \code{hs} arm\\
11 & \code{hsxptari} & rad & true $x'$, \code{hs} arm\\
12 & \code{hsytari} & cm & $z_t\sin\theta^0$ of the \code{hs} arm\\
13 & \code{ssdelta} & \% & reconstructed $\delta$, \code{ss} arm\\
14 & \code{ssyptar} & rad & reconstructed $y'$, \code{ss} arm\\
15 & \code{ssxptar} & rad & reconstructed $x'$, \code{ss} arm\\
16 & \code{ssytar} & cm & reconstructed horizontal target position, \code{ss} arm\\
17 & \code{ssxfp} & cm & focal-plane $x$, \code{ss} arm\\
18 & \code{ssxpfp} & rad & focal-plane $x'$, \code{ss} arm\\
19 & \code{ssyfp} & cm & focal-plane $y$, \code{ss} arm\\
20 & \code{ssypfp} & rad & focal-plane $y'$, \code{ss} arm\\
21 & \code{ssdeltai} & \% & true $\delta$ after radiation, \code{ss} arm\\
22 & \code{ssyptari} & rad & true $y'$, \code{ss} arm\\
23 & \code{ssxptari} & rad & true $x'$, \code{ss} arm\\
24 & \code{ssytari} & cm & $-z_t\sin\theta^0$ of the \code{ss} arm\\
25 & \code{q} & GeV/$c$ & $|\vec q^\mathrm{rec}|$\\
26 & \code{nu} & GeV & $\nu^\mathrm{rec}$\\
27 & \code{Q2} & (GeV/$c$)$^2$ & $(Q^\mathrm{rec})^2$\\
28 & \code{W} & GeV & $W^\mathrm{rec}$, computed with the proton mass\\
29 & \code{xbj} & -- & $x^\mathrm{rec}$\\
30 & \code{epsilon} & -- & $\varepsilon^\mathrm{rec}$\\
31 & \code{Em} & GeV & $E_m^\mathrm{rec}=\nu^\mathrm{rec}+M_N-E_h^\mathrm{rec}$\\
32 & \code{Pm} & GeV/$c$ & $P_m^\mathrm{rec}=|\vec P_m^\mathrm{rec}|$\\
33 & \code{qi} & GeV/$c$ & vertex $|\vec q|$\\
34 & \code{nui} & GeV & vertex $\nu$\\
35 & \code{Q2i} & (GeV/$c$)$^2$ & vertex $Q^2$\\
36 & \code{Wi} & GeV & vertex $W$, computed with the struck-nucleon mass $M_N$ (no Fermi motion)\\
37 & \code{xbji} & -- & vertex $x$\\
38 & \code{epsiloni} & -- & vertex $\varepsilon$\\
39 & \code{Emi} & GeV & vertex $E_\mathrm{miss}=\nu+M_A-E_h$\\
40 & \code{Pmi} & GeV/$c$ & vertex $P_\mathrm{miss}=|\vec p_h-\vec q|$\\
41 & \code{missmass} & GeV & reconstructed missing mass on the struck nucleon at rest (\secr{sec:recon})\\
42 & \code{mmnuc} & GeV & reconstructed missing mass on the target nucleus (\secr{sec:recon})\\
43 & \code{phad} & GeV/$c$ & $P_h^\mathrm{rec}$\\
44 & \code{t} & GeV$^2$ & reconstructed $-t$ (\secr{sec:recon})\\
45 & \code{ti} & GeV$^2$ & vertex $-t$\\
46 & \code{z} & -- & $E_h^\mathrm{rec}/\nu^\mathrm{rec}$\\
47 & \code{zi} & -- & $E_h/\nu$ at the vertex\\
48 & \code{pt2} & GeV$^2$ & $(P_h^\mathrm{rec})^2\sin^2\theta_{pq}^\mathrm{rec}$\\
49 & \code{pt2i} & GeV$^2$ & $P_h^2\sin^2\theta_{pq}$ at the vertex\\
50 & \code{thetapq} & rad & $\theta_{pq}^\mathrm{rec}$\\
51 & \code{phipq} & rad & $\varphi_{pq}^\mathrm{rec}$\\
52 & \code{thetapqi} & rad & vertex $\theta_{pq}$\\
53 & \code{phipqi} & rad & vertex $\varphi_{pq}$\\
54 & \code{pmpar} & GeV/$c$ & $P^\mathrm{rec}_{m,\parallel}$\\
55 & \code{pmper} & GeV/$c$ & $P^\mathrm{rec}_{m,\perp}$\\
56 & \code{pmoop} & GeV/$c$ & $P^\mathrm{rec}_{m,\mathrm{oop}}$\\
57 & \code{fry} & cm & minus the vertical raster coordinate\\
58 & \code{radphot} & GeV & $E_{\gamma1}+E_{\gamma2}+E_{\gamma3}$\\
59 & \code{pfermi} & GeV/$c$ & $p_F$ with the sign of $\hat p_F\cdot\hat q$ (0 without Fermi motion)\\
60 & \code{siglab} & $\mu$b/(MeV sr$^2$) & $\sigma$ at the vertex, with the Coulomb factor (\secr{sec:xsec})\\
61 & \code{sigcm} & $\mu$b/(MeV$^2$ rad) & the centre-of-mass value written by the model (\secr{sec:xsec})\\
62 & \code{Weight} & see \secr{sec:xsec} & the event weight $w$\\
63 & \code{decdist} & cm & path flown by the hadron before its decay, or the whole path when it did not decay (\secr{sec:rad})\\
64 & \code{Mhadron} & MeV & mass of the particle that reached the detectors (the muon mass after a pion decay)\\
\end{simctable}

For this reaction \code{phad} is the proton momentum, \code{missmass} the missing (pion) mass of \eqr{eq:recmm},
\code{t} and \code{ti} the proton momentum transfer, \code{siglab} the value of \eqr{eq:peedelta} (no cross section),
\code{sigcm} the Blok fit of \eqr{eq:sigblok} (meaningless numbers in the compat mode), \code{decdist} 0 and
\code{Mhadron} $M_p$.

% =====================================================================================================================
\section{Validation in SIMC++}\lbl{sec:validation}
\begin{simctable}{footnotesize}{L{0.46\linewidth}L{0.2\linewidth}L{0.34\linewidth}}
\caption{CTest tests that cover (e,e\pr p)$\pi$ (one test per row).}\lbl{tab:tests}\\
\toprule\rowcolor{simchead}\simcth{test} & \simcth{labels} & \simcth{what it checks}\\\midrule\endfirsthead
\toprule\rowcolor{simchead}\simcth{test} & \simcth{labels} & \simcth{what it checks}\\\midrule\endhead
\bottomrule\endlastfoot
\code{l1\_peedelta} & l1 & \eqr{eq:peedelta} and \code{sigcm} against the Fortran table\\
\code{l1\_transform\_to\_cm\_delta} & l1 & the centre-of-mass transform with the proton\\
\code{l1\_delta\_state} & l1 & the set-up of the Delta runs\\
\code{l1\_generate\_delta\_fixedrng} & l1 & \secr{sec:gen} with a fixed random stream\\
\code{l1\_complete\_ev\_delta} & l1 & \eqr{eq:roots}\\
\code{l1\_complete\_recon\_ev\_delta} & l1 & \eqr{eq:recmm}\\
\code{l1\_complete\_main\_delta} & l1 & \eqr{eq:weight}\\
\code{l1\_replay\_event\_capture\_delta\_d1} & l1, replay & every routine of captured events\\
\code{l1\_simcrun\_replay\_delta\_d1} & l1, replay & the complete run on a captured random stream, bitwise\\
\code{l1\_montecarlo\_replay\_delta\_d1} & l1, replay & the spectrometer calls of the capture\\
\code{l1\_simcrun\_replay\_delta\_d2} & l1, replay & the same, second kinematics\\
\code{l1\_simcrun\_ranlux\_up\_delta\_h\_d1\_rad\_A} & l1, replay, slow, upstream\_layout & full run, RANLUX, compat, against the upstream Fortran run\\
\code{regression\_identity\_delta\_h\_d1\_rad} & regression & compat-mode outputs against the baseline\\
\code{regression\_corrected\_delta\_h\_d1\_rad} & regression & default-mode outputs against the baseline\\
\code{regression\_upstream\_u59\_delta\_d\_d1\_rad\_A} & regression\_upstream, slow & against the upstream run\\
\code{regression\_shipped\_shipped\_test\_delta\_pct} & regression\_shipped, slow & a shipped Delta input\\
\code{fortran\_diff\_D048\_F1\_delta\_DA\_stop} & fortran\_diff, unit & D, He-3, C-12 runs: energy--momentum balance of every event\\
\code{fortran\_diff\_D048\_F2\_delta\_tail1\_egamma} & fortran\_diff, unit & the tail-1 limits of \eqr{eq:tail1}\\
\code{fortran\_diff\_D048\_F3\_delta\_sigblok\_args} & fortran\_diff, unit & the arguments of \eqr{eq:sigblok}\\
\code{fortran\_diff\_D048\_F9\_central\_sigcc\_delta} & fortran\_diff, unit & the central value\\
\end{simctable}
Recorded in the validation ledger (\code{project\_control/VALIDATION\_LEDGER.md}, entry of 2026-10-17, task P11.1): the
seven L1 tables bitwise in the compat mode (about $1.7\times10^6$ values), both captures replayed, four L2 comparisons
with the Fortran passing in the compat mode; in the default mode the tail-1 correction multiplies \code{wtcontr} of
the radiative runs by 1.76 and 1.84, the other corrections are deterministic (\code{sigcm}, central value) or at the
0.5\,\% level (hadron $\delta$ cut).

% =====================================================================================================================
\section{Fortran origin and Fortran findings}\lbl{sec:origin}
\begin{simctable}{footnotesize}{L{0.42\linewidth}L{0.58\linewidth}}
\caption{Fortran routine $\to$ C++ function (one routine per row; the routines shared with the pion reaction are in
note 04, the common ones in \ctabr{tab:map}).}\lbl{tab:map}\\
\toprule\rowcolor{simchead}\simcth{Fortran (upstream 59ccdf6)} & \simcth{C++ (src/...)}\\\midrule\endfirsthead
\toprule\rowcolor{simchead}\simcth{Fortran (upstream 59ccdf6)} & \simcth{C++ (src/...)}\\\midrule\endhead
\bottomrule\endlastfoot
\file{dbase.f:182}-190 (flags, warning) & \csrc{core/RunSetup.cpp: DeriveReactionFlags}; \csrc{setup/SetupBuilder.cpp: BuildRunSetup}\\
\file{dbase.f:301}-304 (masses) & \csrc{setup/SetupReaction.cpp: BuildReactionAndKinematics}\\
\file{dbase.f:787}-796 (announcement, stop) & \csrc{setup/SetupBuilder.cpp: BuildRunSetup}\\
\file{event.f} generate (delta flags) & \csrc{reactions/pion/PionProduction.cpp: PionProduction::generate}\\
\file{event.f} complete\_ev, l.682-684 & \csrc{reactions/pion/PionProduction.cpp: PionProduction::completeEv}\\
\file{event.f} complete\_main, l.1514-1516 & \csrc{reactions/delta/DeltaProduction.cpp: DeltaProduction::completeMain}\\
\file{physics\_delta.f} peedelta & \csrc{physics/delta/DeltaModel.cpp: peedelta}; plug-in \csrc{physics/delta/DeltaModel.cpp: DeltaModel::sigccMain}\\
\file{physics\_pion.f} sig\_blok & \csrc{physics/pion/PionCrossSections.cpp: sigBlok}\\
\file{jacobians.f} transform\_to\_cm & \csrc{reactions/EventChain.cpp: transformToCm}\\
\file{radc.f:288}-303 (tail 1) & \csrc{radiation/Radc.cpp: generateRad}\\
\end{simctable}

\begin{simcfindings}{footnotesize}{L{0.18\linewidth}L{0.5\linewidth}L{0.32\linewidth}}
\caption{Fortran findings that touch (e,e\pr p)$\pi$ (one finding per row; the common ones are in \ctabr{tab:findings}).}\lbl{tab:findings}\\
\toprule\rowcolor{simchead}\simcth{id} & \simcth{Fortran behaviour} & \simcth{default / compat}\\\midrule\endfirsthead
\toprule\rowcolor{simchead}\simcth{id} & \simcth{Fortran behaviour} & \simcth{default / compat}\\\midrule\endhead
\bottomrule\endlastfoot
\finding{D-048 F4} & no cross section in the weight & kept in both (open question)\\
\finding{D-048 F1 Delta D/A STOP} & every D or $A\ge3$ run stops (pion flags tested) & runs / stop\\
\finding{D-048 F1 Delta D/A data} & momentum distribution and spectral function not read for the Delta & read / not read\\
\finding{D-048 F2} & no tail-1 limits: every tail-1 try fails & pion limits / Fortran\\
\finding{D-048 F3} & \code{sig\_blok} arguments of the wrong kind: \code{sigcm} meaningless & intended values / Fortran\\
\finding{D-048 F5} & random root choice without a factor 2 & kept in both\\
\finding{D-048 F5b} & a tail-1 photon draws a new root & kept in both\\
\finding{D-048 F6} & quote debris in the D/A warning & clean text / Fortran\\
\finding{D-048 F7} & recoil mass of the charged pion although the missing pion is a $\pi^0$ & kept in both\\
\finding{D-048 F8} & \code{which\_pion} read by the Delta reaction & kept in both\\
\finding{D-048 F9 CENTRAL.sigcc} & central value with $\varepsilon=-t=\varphi_{pq}=0$ & central kinematics / Fortran\\
\finding{D-048 F9 stale fields} & stale tree-unrelated fields; $n_\mathrm{tgt}=1$ also for D/A & the coded weight\\
\finding{D-048 F11} & the stop exits with status 0 & no stop / exit status 3\\
\finding{P11.1 tfcos} & a comparison in quadruple precision & reproduced\\
\finding{D-006 pass\_cuts} & hadron $\delta$ upper cut with the electron limit & hadron limit / Fortran\\
\end{simcfindings}

\begin{simcrefs}
\simcref{brauel} P. Brauel et al., Z.\ Phys.\ C \textbf{3}, 101 (1979) (cited at \fsrc{physics\_pion.f:408}).
\end{simcrefs}
\end{document}
